\documentclass{article}[notitlepage]

\usepackage{amsmath}
\usepackage{amssymb}
\usepackage{amsthm}
\usepackage{derivative}
\usepackage{kotex}
\usepackage{fullpage}
\usepackage{mdframed}

\newtheorem*{theorem}{Theorem}
\theoremstyle{definition}
\newtheorem*{definition}{Definition}

\mdfdefinestyle{MyFrame}{
    outerlinewidth=2pt,
    roundcorner=20pt
}
\def \dd{\mathrm d}

\title{Renewal Process for Wildfire Spreading}
\author{이동하}
\date{}

\begin{document}
\maketitle
The occurance of wildfire, which is modeled with the probability, can be viewed as a kind of point processes, especially the renewal process. The renewal process refers to the point process whose inter-point time interval distribution is not i.i.d.(independent and identically distributed).

\section{Theory}

\subsection{Renewal Process}
Consider two different incidents of a point process which is occured at different time points $t_0$ and $t$, where $t_0 < t$. Define the followings:
\begin{enumerate}
    \item \textbf{next point distribution $p(t|t_0)$:} the probability distribution for time point of the next event, given event at $t_0$,
    \item \textbf{survivor function $S(t|t_0)$:} the probability of not-occurance of event until time $t$, given event at $t_0$, and
    \item \textbf{hazard function $\rho(t|t_0)$:} the probability distribution for the chance of occurance of event at time point $t$, given event at $t_0$.
\end{enumerate}
Observe that the following equations are plausible:
\[p(t|t_0) = S(t|t_0) \rho(t|t_0), \qquad S(t|t_0) = 1 - \int^{t}_{t_0} p(\tau|t_0) \dd \tau.\]
From these two equations, the differential equation
\[\pdv*{S(t|t_0) }{t} = - S(t|t_0) \rho(t|t_0)\]
can be established. Since $S(t_0|t_0) = 1$, solving the differential equation, we have
\[S(t|t_0) = \exp\left[-\int^t_{t_0} \rho(\tau|t_0) \dd{\tau}\right],\]
and hence
\[p(t|t_0) = \rho(t|t_0) \exp\left[-\int^t_{t_0} \rho(\tau|t_0) \dd{\tau}\right].\]
Note that if $\rho(t|t_0)$ does not depend on $t_0$, that renewal process is Poisson process.


\end{document}